% ED26-BM-RAT2021-10-1, exact MIT 2021 chapter, printed pp. 579--590.
\chapter{Social Epistemology}
\label{sec:p18-rationality-10-1}

A group can reach unanimous agreement without gaining evidence, and it
can possess useful information without forming a coherent collective
judgment. Franz Dietrich and Kai Spiekermann's chapter~10.1 distinguishes
these problems rather than treating agreement as knowledge. Its central
subject is the formation and truth-tracking performance of group beliefs
and other correctness-apt collective decisions. Testimony, disagreement,
and social justification enter that account, but the chapter expressly
does not provide a comprehensive theory of individual testimonial
justification or of rational peer-disagreement revision
\parencite[pp.~579--580 and n.~1, p.~586]{DietrichSpiekermann2021Social}.

The reconstruction below follows the complete numbered 2021 chapter,
including its nine notes and references. The formal demonstrations are
conditional mathematical results; the examples of institutions are
mechanisms and qualified applications; the selected empirical results
retain their populations, tasks, measures, and inferential limitations.
Neither a voting theorem nor a Bayesian network is the whole subject.

\section{Whose knowledge, and what counts as success?}
\label{sec:p18-r101-scope}

\begin{hbdefinition}{def:p18-r101-holders}{Knowledge in groups and knowledge of groups}
\emph{Epistemology in groups} concerns an individual's knowledge or
belief as affected by social inputs: testimony, another person's
disagreement, or an opinion spreading through a social network.
\emph{Epistemology of groups} concerns beliefs or knowledge attributed
to the group itself. The latter attribution has different senses:
\emph{aggregate beliefs} summarize members' attitudes;
\emph{common beliefs} involve members' beliefs together with their
common awareness; and \emph{corporate beliefs} belong literally to a
group construed as an agent. A committee, firm, or state might qualify
for the last attribution; a survey sample or a prediction market need
not. Group agency requires an additional account, not merely a large
number of people agreeing.
\end{hbdefinition}

The authors use belief-revision models to mark the formal problem of
learning within groups, and aggregation models to mark the formal
problem of forming group attitudes. Their review concentrates on the
second. The distinction does not prohibit interactions between them:
individuals may revise their beliefs during discussion before a committee
aggregates its final judgments. Nor does the aggregate/common/corporate
distinction imply that a market literally knows its forecast. The chapter
attributes the three-way distinction to List and sets aside the disputed
metaphysics of group knowledge
\parencite[p.~579]{DietrichSpiekermann2021Social}.

\begin{hbdefinition}{def:p18-r101-correctness}{Narrow and broad epistemic decisions}
Narrow social epistemology assesses social beliefs. Broad social
epistemology also assesses collective actions and, where appropriate,
other attitudes. For an action or attitude to enter the broad account,
there must be a relevant sense in which it is correct or incorrect.
The correctness standard must be sufficiently stable and independent
of the very procedure being evaluated. Logical, empirical, and normative
facts can supply different kinds of standard. Calling a decision
``social'' does not establish voluntary control over the resulting belief.
\end{hbdefinition}

A court's belief that a defendant committed an act differs from its
decision to impose a sentence. A scientific community's acceptance of a
law differs from its choice of the next experiment. The authors' examples
range from mathematical belief, through a central bank's inflation
forecast, to doctors reallocating resources from research to treatment.
The last case involves empirical and normative considerations in addition
to logical ones. Their broad conception allows such questions without
settling whether normative facts are objective, constructed, universal,
or culturally relative. A constructed standard must precede the assessed
procedure; if the procedure itself creates the standard, its outcome is
not tracking an independently specified truth. Note~4 also permits a
position that denies normative facts while assessing normative judgments
as correct or incorrect acts. Note~3 raises, without resolving, comparable
questions about intentions and other judgment-sensitive attitudes
\parencite[pp.~580 and 586--587, nn.~2--4]{DietrichSpiekermann2021Social}.

\begin{hbwarning}{def:p18-r101-three-standards}{Coherence, truth, and legitimacy are different standards}
In this chapter \emph{formal rationality} means the relevant internal
coherence: consistency and suitable closure for binary judgments,
probability coherence for credences, or transitivity and possibly
completeness for rankings. Coherence does not establish that the beliefs
are true or the rankings correct. Individual coherence does not guarantee
collective coherence. Epistemic performance, in turn, does not by itself
settle democratic legitimacy or procedural fairness.
\end{hbwarning}

The authors treat formal rationality as a necessary step toward a
correct collective attitude of the appropriate type, not as a sufficient
condition. A consistent but false theory illustrates the gap. The
relationship can also be indirect: a personally suboptimal revision policy
may preserve useful diversity for a collective search. Social knowledge
inherits disputes about justified true belief, Gettier-type coincidence,
evidence, and reliable procedures. Reliability concerns how an attitude
is produced across relevant cases; justification by evidence concerns the
grounds available for it. A correct vote by accident settles neither
issue. The chapter does not claim to resolve evidentialism versus
reliabilism or to give a sufficient definition of social knowledge
\parencite[p.~580 and n.~5, p.~587]{DietrichSpiekermann2021Social}.

The democratic comparison makes the difference concrete. An epistemic
justification values procedures because their outcomes track an
independent standard; a proceduralist justification instead appeals to
such considerations as equal participation and fair treatment of options.
The chapter's hypothetical infallible oracle would be excellent at the
first task and wholly undemocratic at the second. Its performance would
not prove a right to govern. Hybrid positions remain possible, and an
empirical finding about a procedure needs a separate normative argument
before becoming a claim about legitimacy
\parencite[pp.~580--581]{DietrichSpiekermann2021Social}.

\section{Explicit aggregation: inputs, outputs, and their conflicts}
\label{sec:p18-r101-aggregation}

\begin{hbdefinition}{def:p18-r101-binary-rules}{Binary belief aggregation}
Members answer yes or no to a proposition. The collective output may
also abstain. Majority rule adopts the more numerous answer and abstains
on a tie. An asymmetric supermajority rule requires a specified greater
than half threshold for yes and otherwise returns no. A symmetric
supermajority rule requires the same threshold for either answer and
otherwise abstains. Unanimity is the corresponding extreme threshold.
The chapter simplifies by disallowing individual abstention; that is a
model restriction, not a statement that actual discussants cannot suspend
judgment.
\end{hbdefinition}

These rules expose a tradeoff between accepting a claim, rejecting it,
and withholding a collective answer. A symmetric high threshold increases
abstention, whereas an asymmetric threshold privileges one answer. The
choice is not justified merely by observing that a jury or economic
panel has an answer to aggregate. Its consequences depend on accuracy,
dependence, and the costs or purposes of the collective decision.

For several propositions, separate tallies can violate their logical
connections. In the chapter's court illustration, let \(p\) say that
the defendant committed an act, \(q\) that such an act breaches the law,
and \(r\) that the defendant breached the law. To make the precise
premise-based example explicit, use the declared agenda constraint
\(r\leftrightarrow(p\land q)\). This constraint is part of the
model, not a universal claim about legal responsibility. Majority
acceptance of \(p\) and \(q\) together with rejection of \(r\) is
then inconsistent. Withholding an answer on \(r\) instead would leave
the set deductively unclosed relative to this agenda.

\begin{hbtheorem}{thm:p18-r101-aggregation-failures}{proved=annex:p18.r101.aggregation}{Three distinct failures of simple aggregation}
There are three-member profiles satisfying
\(r\leftrightarrow(p\land q)\) at every member but violating it under
propositionwise majority. There are three-member strict rankings whose
pairwise-majority relation cycles. There is a plurality profile in which
splitting support among similar alternatives causes a less-supported
unrelated option to win. These are different counterexamples, not one
general impossibility theorem.
\end{hbtheorem}

For the judgment example, use the individual rows
\((p,q,r)=(1,1,1),(1,0,0),(0,1,0)\). The majority row is
\((1,1,0)\). A premise-based rule first aggregates \(p\) and \(q\),
then infers \(r\); here it accepts \(r\) against the majority on
that conclusion. It restores this agenda's consistency by a justified
choice of priority, not by respecting every issuewise majority.
The source also names distance-based, sequential, relevance-based, and
approximate-majoritarian approaches. Their differences concern how
interconnected issues are handled; none is supplied merely by the name
``holistic.'' The chapter's rough references to general judgment-aggregation
impossibility results are not the assertion that every possible group
rule is defective. Their agenda and axiom hypotheses matter, and the
finite counterexample does not prove their general characterization
\parencite[pp.~581--582]{DietrichSpiekermann2021Social}.

\subsection{Probabilistic pooling: two different coherence demands}

Binary judgment aggregation and probabilistic pooling must not be
conflated. Local averaging can preserve probability coherence even where
local majority fails to preserve logical consistency.

\begin{hbdefinition}{def:p18-r101-pools}{Finite linear and geometric pools}
Let \(\Omega\) be a finite nonempty set of worlds. Member \(i\) has
a probability mass function \(P_i\), and fixed weights \(w_i\geq0\)
sum to one. The linear pool is
\[
 L(A)=\sum_i w_iP_i(A)\qquad(A\subseteq\Omega).
\]
When every \(P_i(\omega)>0\), define the geometric pool by
\[
 G(\omega)=\frac{\prod_iP_i(\omega)^{w_i}}
 {\sum_{\nu\in\Omega}\prod_iP_i(\nu)^{w_i}},
 \qquad G(A)=\sum_{\omega\in A}G(\omega).
\]
Strict positivity here avoids undefined normalization or conflicting
zero-probability supports. Pooling is a rule for forming a collective
distribution, not automatically an individual's update from independent
reports.
\end{hbdefinition}

\begin{hbtheorem}{thm:p18-r101-pooling}{proved=annex:p18.r101.pooling}{Coherence, common conditioning, and partition dependence}
The finite linear pool is a probability function and respects coarsening
by summation. With fixed weights it need not commute with conditioning
all members on the same positive-probability event. The strictly positive
finite geometric pool does commute with such common conditioning, but
pooling after refining the worlds need not give the same coarse-event
probabilities as pooling on the coarse space.
\end{hbtheorem}

For equal weights and worlds \(a,b,c\), take
\(P_1=(0.6,0.3,0.1)\), \(P_2=(0.2,0.3,0.5)\), and
\(E=\{a,b\}\). Conditioning their linear pool gives
\(L(a\mid E)=4/7\); pooling their conditionals gives \(8/15\).
The difference concerns the order of two operations. It does not show
that an averaged distribution violates the probability axioms.
For partition dependence, take both agents' coarse probabilities of
\(A,B\) to be \((1/2,1/2)\), then split \(A\) into \(a_1,a_2\)
with refined distributions \((0.45,0.05,0.5)\) and
\((0.05,0.45,0.5)\). The coarse geometric pool assigns \(1/2\)
to \(A\); the refined pool assigns \(3/8\).

The chapter contrasts probabilism, a constraint at a time, with
conditionalization, a constraint on change. Its cited axiomatic
characterizations of linear pooling and of suitably well-behaved Bayesian
groups have additional hypotheses. The finite commutation result is not
a proof that geometric pooling is the uniquely rational policy in every
setting, especially not when members learn different evidence or change
their weights. This interface connects to
Definition~\ref{def:p15-rat41-probabilism} and the individual dynamics
discussion in Chapter~\ref{sec:p15-rationality-4-1}; pooling different
people and revising one person's belief are different tasks
\parencite[p.~582]{DietrichSpiekermann2021Social}.

\subsection{Evaluations and direct choices}

The source distinguishes ordinal evaluations, represented by rankings,
from absolute evaluations, represented by grades in a common ordered
scale. Wine quality, moral assessment, size, and candidates illustrate
different possible criteria; treating an evaluation as epistemic assumes
that some relevant standard of correctness exists. Pairwise majority of
ordinal rankings can cycle: rankings \(a>b>c\), \(b>c>a\), and
\(c>a>b\) collectively favor \(a\) over \(b\), \(b\) over \(c\),
and \(c\) over \(a\). Arrow-type restrictions are source-cited
general results, not the unqualified claim that every aggregation
procedure fails. Holistic procedures change issuewise independence;
absolute grading changes the informational representation.

The chapter's median-grade illustration assigns a wine the grades
exquisite, exquisite, and drinkable; the social grade is exquisite.
It presents this as the basic median step of majority judgment. It is
not a complete specification of all tie-breaking refinements of that
voting system. Likewise, its two-grade approval illustration awards
approval to an option with a maximal number of approvals; it does not
imply that an option approved by someone must be a winner. A common
grade vocabulary and its comparable interpretation are substantive
inputs, not supplied by ordinal rankings alone
\parencite[pp.~582--583]{DietrichSpiekermann2021Social}.

Direct social-choice rules output a winning alternative or set, rather
than a complete social ranking. Plurality lets each person choose one
option and selects a highest-count option. In the theorem's illustrative
profile, forty voters support \(b\) while sixty support a correct
option \(a\); replacing that option by three similar variants with
twenty votes each lets \(b\) win. Whether variants are equally correct
or merely approximate must be declared separately. Approval permits
support for several variants and can avoid that particular vote-splitting
mechanism; it is not an unconditional accuracy guarantee. Ranking-based
choice can face a different problem, absence of a Condorcet winner.
The chapter's note~6 distinguishes these direct action-producing rules
from first forming beliefs and then acting on them
\parencite[p.~583 and n.~6, p.~587]{DietrichSpiekermann2021Social}.

\section{Jury theorems and the dependence that matters}
\label{sec:p18-r101-juries}

\begin{hbdefinition}{def:p18-r101-jury-model}{Correct-vote probability and growing reliability}
For a binary question with exactly one correct answer, \(X_i\) is one
when member \(i\)'s answer is correct and zero otherwise. For an odd
group of size \(n\), majority is correct when
\(S_n=\sum_{i=1}^{n}X_i>n/2\). Write \(M_n(t)\) for this probability
when the indicators are independent with common success probability
\(t\). \emph{Growing reliability} compares this performance as group
size increases; \emph{asymptotic infallibility} asserts convergence to
one. The second does not mean a finite group cannot err.
\end{hbdefinition}

\begin{hbtheorem}{thm:p18-r101-jury}{proved=annex:p18.r101.jury}{Classical majority reliability and a finite common-cause replacement}
For \(1/2<t<1\),
\[
 M_n(t)=\sum_{k=(n+1)/2}^{n}{n\choose k}t^k(1-t)^{n-k}
\]
strictly increases through odd group sizes and tends to one.
Now let a finite-valued problem variable \(\Pi\) have probabilities
\(\mu_j\). Conditional on \(\Pi=j\), suppose votes are independent
with common competence \(t_j\). Their unconditional majority reliability
is \(\sum_j\mu_jM_n(t_j)\), with limit
\[
 \sum_{t_j>1/2}\mu_j+\tfrac12\sum_{t_j=1/2}\mu_j.
\]
If for every reflected pair \(1/2+d,1/2-d\) the probability mass at
the upper value is at least that at the lower, this reliability weakly
increases through odd sizes. Average individual competence above one
half alone does not establish that conclusion.
\end{hbtheorem}

The assumptions determine what the result says. Majority correct-vote
independence is stochastic independence relative to the assessed state
or problem, not merely the absence of conversation. Shared evidence,
background assumptions, a common opinion leader, or even a common
noninformational influence can correlate errors. Adding more readers
of one misleading report does not create new evidence.

Dietrich and Spiekermann's 2013 treatment replaces independence
conditional only on the correct state by independence conditional on a
problem variable that also represents relevant common circumstances.
Its competence requirement uses a reflected-distribution dominance
around one half, not just a high mean. The finite theorem above makes
this actual mechanism explicit. The original permits more general
problem variables and distributions; that general treatment is not
silently replaced by a claim that all actual social influence can be
encoded in the finite model. Its causal-network discussion motivates
screening off; conditional independence still needs substantive
justification. Direct interpersonal influence or unrepresented common
causes can defeat it
\parencite[pp.~94--103]{DietrichSpiekermann2013Defensible}.

For a concrete fallible-but-improving case, let competence be \(0.8\)
on problems occurring with probability \(0.8\), and \(0.2\) on the
remaining problems. Individual accuracy is \(0.68\); the limiting
majority accuracy is \(0.8\), not one. For the source's counterexample
to a mere mean-competence premise, give equal probability to competence
\(0.99\) and \(0.49\). Individual accuracy is \(0.74\), but the
majority limit is \(0.5\). An average skill parameter is therefore
not an adequate substitute for the distributional premise.

The 2021 chapter distinguishes the two reliability hypotheses, criticizes
the interpretation of infallibility as a realistic empirical promise,
and observes that heterogeneous competence alone does not repair the
dependence problem. It notes extensions to plurality, approval,
distance-based and scoring rules. The existence of some jury theorem
for a rule is not evidence that its assumptions hold for a particular
committee. Note~7 further distinguishes assessing a fixed rule from
searching for an epistemically optimal rule with differential
competences. Correct outcomes, sincere reports, fairness, and incentives
are separate properties
\parencite[pp.~583--584 and n.~7, p.~587]{DietrichSpiekermann2021Social}.

\section{Institutions that do more than count judgments}
\label{sec:p18-r101-institutions}

\subsection{Deliberation and its empirical boundary}

Deliberation exchanges reasons before or instead of aggregation. The
chapter identifies five potential contributions: making private grounds
available; exposing bad reasons and inconsistencies; reducing biases or
opinion-leader effects; changing the agenda by discovering or reframing
options; and producing agreement about the structure of the decision
space without agreement on the final answer. These are distinct
mechanisms. Discussion can instead create dependence or suppress useful
information, so their availability is not a theorem of epistemic success.
The source cites deliberative-polling research as encouraging while also
retaining polarization objections
\parencite[p.~584]{DietrichSpiekermann2021Social}.

The selected field example is Esterling, Neblo, and Lazer's 2006
immigration-policy experiment: twelve US House members held twenty
online sessions with constituent groups of eight to thirty. Among
2,222 Internet-panel constituents, assignment offered background
information plus a 35-minute representative discussion and 25-minute
peer chat, information alone, or neither. Six factual items measured
policy knowledge one week later. A principal-stratification model
addressed noncompliance and nonresponse under conditional
missing-at-random assumptions. Estimated knowledge gain versus
information alone was 0.36 standard deviations (reported 95\% interval
0.11--0.61); versus controls it was 0.61 (0.38--0.84).
The treatment bundles access, information, and discussion; its mechanisms
are not separately randomized. The inference is short-term and restricted
to Internet-connected constituents in the selected districts, not all
citizens or all deliberation. The chapter itself warns that deliberative
process and opportunities for influence are difficult to disentangle
\parencite[pp.~485--495, 497--500]{EsterlingNebloLazer2011Informed}
\parencite[p.~584]{DietrichSpiekermann2021Social}.

\subsection{Distributed search}

Distributed search divides epistemic labor instead of inducing everyone
to investigate the same possibility. The source emphasizes large search
spaces, expensive or lengthy tests, and useful by-products of work in
different locations. Scientific research and the historical federalist
``laboratory'' image are illustrations, not equivalent institutions or
a controlled experiment comparing political systems. A reward structure
that makes chasing fashionable, well-funded projects individually
attractive may leave valuable alternatives unexplored. Diversity,
replication, and longer-horizon work can therefore be collective design
considerations even when each researcher responds sensibly to existing
incentives. The chapter cites models by Kitcher, Strevens, Hong and Page,
Weisberg and Muldoon, and Zollman as different ways of investigating this
tension. A result from one search landscape does not establish universal
benefits of maximum diversity, nor does temporary disagreement necessarily
mean epistemic failure
\parencite[pp.~584--585]{DietrichSpiekermann2021Social}.

\subsection{Prediction markets}

Prediction markets aggregate through buying and selling event-contingent
contracts, rather than through explicit yes/no reports. Profits and losses
can encourage acquiring information and correcting prices. Market prices
may represent implicit collective forecasts without making the market a
corporate agent. Their interpretation as probabilities depends on market
and behavioral assumptions; trading is not identical to truthful
testimony. The chapter's claim that counterjudgmental trading normally
loses money is an incentive rationale, not a proof excluding manipulation,
wealth effects, or illiquid markets
\parencite[p.~585]{DietrichSpiekermann2021Social}.

Dreber and colleagues used two-week markets in 2012 and 2014 for
44 scheduled psychology replications. Researchers recruited through
Open Science Framework and replication-project lists received \$100
trading endowments; 47 and 45 actively traded. Success generally meant
\(p<0.05\) in the original direction, with one original null-result
exception. Among 41 completed replications, market predictions at a
50\% cutoff were correct for 29 (71\%; binomial test \(p=0.012\)).
For 40 comparable cases, absolute error was lower than the premarket
survey's (paired \(t=-2.558\), \(p=0.015\)). These are selected
psychology tasks, not a population accuracy guarantee. Three outcomes
were unfinished; one target changed; trader eligibility and disclosure
restrictions mattered. Forecasting replication is not establishing
truth of a hypothesis, and the study does not isolate incentives from
information exchange or aggregation
\parencite[pp.~15343--15346]{DreberEtAl2015ReplicationMarkets}.

The chapter mentions the 2016 Brexit and US-election forecasts as
warnings against infallibility. Its note~9 correctly qualifies the
language of ``wrong prediction'': a low probability for an event is
compatible with that event occurring. Assessing systematic
underestimation requires a forecast-performance standard and suitable
data, not one surprise alone. Its cited Iowa-market website and historical
event illustrations are not newly analyzed forecast series here
\parencite[p.~585 and n.~9, p.~587]{DietrichSpiekermann2021Social}.

\section{Failures of reports, reasons, and social interpretation}
\label{sec:p18-r101-failures}

Strategic voting is not merely bad belief formation. The chapter explains
that a juror who wants the group to be correct can have a reason to vote
against her own private assessment, because the vote's effect depends
on pivotal circumstances and the decision rule. Shared outcome-oriented
goals do not ensure sincere reporting. This possibility is distinct from
a prisoner's dilemma with conflicting interests. The source does not
declare every such strategy socially harmful: the verdict depends on the
information model and utilities. A separate concern for expressing one's
own judgment can favor sincerity even where purely consequentialist
voting would not. The 2013 jury model explicitly leaves strategic voting
outside its analysis; when votes are not sincere, its indicators can
instead represent correct private judgments, which do not automatically
become the actual votes
\parencite[pp.~585--586]{DietrichSpiekermann2021Social}
\parencite[p.~93]{DietrichSpiekermann2013Defensible}.

Social pressure can affect which reasons are offered. Information already
shared may be easier to discuss than a crucial privately held item;
someone anticipating controversy may remain silent. The chapter uses
hidden-profile and group-polarization literature to identify possible
failure mechanisms. It does not report a single task, population, or
effect estimate that warrants a general law that deliberation invariably
suppresses minority evidence. Silence, repeated reports, a private signal,
and an independent corroboration must therefore be distinguished.

Motivated cognition can operate earlier, influencing which information
is sought and how it is interpreted. The source's climate-risk example
attributes increased cultural divergence to better scientific and numerical
skills. Kahan and colleagues' original article and supplement describe a
January 2010 observational survey of 1,540 US adults recruited through
Knowledge Networks' online panel, eight basic-science items, fourteen
numeracy questions, two six-item worldview scales, and 0--10 climate-risk
ratings. Its OLS model of standardized risk includes a science/numeracy
by hierarchy interaction of \(-0.05\) (\(t=-2.30\)); the corresponding
individualism interaction is \(-0.02\) (\(t=-1.12\)). The two
interactions jointly add explanatory power, but they are not both
individually significant. Missing worldview/risk items were multiply
imputed. These observational associations do not demonstrate that
teaching science causes biased processing or determine any respondent's
actual information-search mechanism. Cultural-worldview measures are
not interchangeable with party labels. The original authors themselves
describe identity-protective processing as a provisional explanation
requiring experimental corroboration, rather than an observed causal
mechanism identified by this survey
\parencite[author postprint, pp.~6--10]{KahanEtAl2012Climate}
\parencite[supplement, pp.~1--12, Tables S1--S4]{KahanEtAl2012Climate}
\parencite[p.~586]{DietrichSpiekermann2021Social}.

\begin{hbdefinition}{def:p18-r101-injustice}{Two epistemic injustices and two refusals of truth-seeking}
The chapter adopts Fricker's distinction between \emph{testimonial
injustice}, where identity prejudice causes a person's testimony to be
discounted, and \emph{hermeneutic injustice}, where inadequate shared
interpretive resources impede understanding or communicating an
experience. It illustrates the latter with experience of sexual
harassment before an adequate public concept was available. These are
conceptual/historical illustrations, not experimental effect sizes.

\emph{Epistemic skepticism}, as used here, doubts that there is an
accessible truth to track; \emph{epistemic nihilism} treats truth as
unimportant even if accessible. Their difference matters: a skeptical
objection offers a disputed epistemic position, whereas indifference to
truth can resist reason-giving by rejecting its aim.
\end{hbdefinition}

Discounting a source solely because of a prejudicial identity judgment
can deprive a group of information; lack of interpretive resources can
prevent the information from becoming an intelligible report. This is
not the prescription to accept every testimony uncritically. The
chapter's final discussion links disregard for truth to Frankfurt's
account of bullshit and tentatively to fake news. Its moral appraisal
and proposed explanatory link remain the authors' judgments; the
chapter supplies no experiment or population estimate showing that
all misinformation is caused by nihilism. Mistake, deception,
skepticism, strategic expression, and indifference remain distinct
\parencite[p.~586]{DietrichSpiekermann2021Social}.

\section{Using the treatment without collapsing its neighboring tasks}
\label{sec:p18-r101-use}

\begin{hbmethod}{def:p18-r101-selection}{Select an epistemic procedure at its actual scope}
Identify the holder and intended output: one person's revised belief,
a collective judgment set, a probability distribution, a ranking, or
an action. Specify the correctness standard separately from procedural
fairness, legitimate authority, and the costs of errors. Record agenda
constraints, possible abstentions, and the information each input
actually supplies. Trace common sources and direct influence before
treating multiple reports as independent. Select a rule with its genuine
tradeoffs: conclusion priorities for premise-based judgment, fixed
weights and support conditions for pooling, grading comparability for
median assessments, or market settlement conditions for forecasts.
For a performance claim, distinguish an assumption-based theorem from
observations on a specified task/population. After deliberation, check
whether the agenda, available evidence, or dependence structure changed.
Preserve unshared or dissenting reasons rather than using agreement alone
as an acceptance criterion.
\end{hbmethod}

Definitions~\ref{def:p18-expert-opinion-scheme} and
\ref{def:p18-witness-testimony-scheme} assess the grounding of a
particular report: domain-relevant expertise, opportunity to observe,
credibility, and contrary evidence are not supplied by the number of
repeaters. This chapter adds the collective issue of what an aggregation
rule does with reports and their connections. It does not replace the
scheme's questions with a vote. Similarly,
Definition~\ref{def:p18-r55-norms} separates procedural, epistemic,
and conventional standards. An orderly discussion can reach an
unsupported conclusion; a coherent probability pool can be poorly
informed; empirically common behavior is not thereby justified.

Peer disagreement is within epistemology in groups, but the selected
chapter does not decide between equal-weight, steadfast, or other
revision norms. Nor does its use of common causes convert common belief
into the formal common knowledge required by an epistemic game model.
Distributed-search benefits do not establish that dissent should always
be retained, and a corporate agent's adopted judgment can differ from
each member's private belief. Those distinctions are part of the
source's boundaries, not missing axioms to be filled by a template.

The useful success case is a group that preserves relevant information,
uses a coherent representation, supplies justified dependence and
competence assumptions, and assesses its output against a separately
identified standard. The corresponding failures are voting on connected
claims without their constraints, counting copies of one source as
independent evidence, mistaking unanimous belief for truth, inferring
legitimacy from predicted accuracy alone, or transferring a measured
benefit to another task or population. The chapter's full argument is
the conditional evaluation of these practices and institutions, not an
unconditional wisdom-of-crowds endorsement.
